% Options for packages loaded elsewhere
\PassOptionsToPackage{unicode}{hyperref}
\PassOptionsToPackage{hyphens}{url}
\documentclass[
  ignorenonframetext,
]{beamer}
\newif\ifbibliography
\usepackage{pgfpages}
\setbeamertemplate{caption}[numbered]
\setbeamertemplate{caption label separator}{: }
\setbeamercolor{caption name}{fg=normal text.fg}
\beamertemplatenavigationsymbolsempty
% remove section numbering
\setbeamertemplate{part page}{
  \centering
  \begin{beamercolorbox}[sep=16pt,center]{part title}
    \usebeamerfont{part title}\insertpart\par
  \end{beamercolorbox}
}
\setbeamertemplate{section page}{
  \centering
  \begin{beamercolorbox}[sep=12pt,center]{section title}
    \usebeamerfont{section title}\insertsection\par
  \end{beamercolorbox}
}
\setbeamertemplate{subsection page}{
  \centering
  \begin{beamercolorbox}[sep=8pt,center]{subsection title}
    \usebeamerfont{subsection title}\insertsubsection\par
  \end{beamercolorbox}
}
% Prevent slide breaks in the middle of a paragraph
\widowpenalties 1 10000
\raggedbottom
\AtBeginPart{
  \frame{\partpage}
}
\AtBeginSection{
  \ifbibliography
  \else
    \frame{\sectionpage}
  \fi
}
\AtBeginSubsection{
  \frame{\subsectionpage}
}
\usepackage{iftex}
\ifPDFTeX
  \usepackage[T1]{fontenc}
  \usepackage[utf8]{inputenc}
  \usepackage{textcomp} % provide euro and other symbols
\else % if luatex or xetex
  \usepackage{unicode-math} % this also loads fontspec
  \defaultfontfeatures{Scale=MatchLowercase}
  \defaultfontfeatures[\rmfamily]{Ligatures=TeX,Scale=1}
\fi
\usepackage{lmodern}
\ifPDFTeX\else
  % xetex/luatex font selection
\fi
% Use upquote if available, for straight quotes in verbatim environments
\IfFileExists{upquote.sty}{\usepackage{upquote}}{}
\IfFileExists{microtype.sty}{% use microtype if available
  \usepackage[]{microtype}
  \UseMicrotypeSet[protrusion]{basicmath} % disable protrusion for tt fonts
}{}
\makeatletter
\@ifundefined{KOMAClassName}{% if non-KOMA class
  \IfFileExists{parskip.sty}{%
    \usepackage{parskip}
  }{% else
    \setlength{\parindent}{0pt}
    \setlength{\parskip}{6pt plus 2pt minus 1pt}}
}{% if KOMA class
  \KOMAoptions{parskip=half}}
\makeatother
\usepackage{longtable,booktabs,array}
\newcounter{none} % for unnumbered tables
\usepackage{calc} % for calculating minipage widths
\usepackage{caption}
% Make caption package work with longtable
\makeatletter
\def\fnum@table{\tablename~\thetable}
\makeatother
\setlength{\emergencystretch}{3em} % prevent overfull lines
\providecommand{\tightlist}{%
  \setlength{\itemsep}{0pt}\setlength{\parskip}{0pt}}
\usepackage{bookmark}
\IfFileExists{xurl.sty}{\usepackage{xurl}}{} % add URL line breaks if available
\urlstyle{same}
\hypersetup{
  hidelinks,
  pdfcreator={LaTeX via pandoc}}

\author{\texorpdfstring{}{}}
\date{}

\begin{document}

\section{Discussion: Cross-Domain Analysis of Cognitive
Integrity}\label{discussion-cross-domain-analysis-of-cognitive-integrity}

\begin{frame}[allowframebreaks]{Discussion: Cross-Domain Analysis of
Cognitive Integrity}
Our cross-domain analysis of ten critical sectors reveals that Goal
Hijacking is not merely a linguistic exploit but a structural corruption
of the OODA Loop \cite{boyd1987patterns}. In every case---from drone
swarms operating at millisecond time scales to diplomatic agents
spanning months of deliberation---the attack vector was a transient
signal that hijacked the agent's \textbf{Orientation} phase, rewriting
its Functional Requirements in real-time. This section synthesizes the
cross-domain findings, identifies universal attack patterns, evaluates
CIF mechanism coverage, and acknowledges limitations.
\end{frame}

\begin{frame}[allowframebreaks]{4.1 Cross-Domain Attack Pattern
Taxonomy}
\protect\phantomsection\label{sec:attack_patterns}
Three universal attack patterns emerge across the ten domains. Each
pattern corresponds to a distinct manipulation of the Axiomatic Design
Matrix \cite{suh2001axiomatic}:

\textbf{Pattern 1: FR Polarity Inversion.} The adversary flips the sign
of a Functional Requirement, transforming a minimization objective into
a maximization objective (or vice versa). The diagonal element
\(A_{ii}\) effectively changes sign. This is the most common pattern,
appearing in five domains.

\textbf{Pattern 2: Constraint Relaxation.} The adversary degrades a hard
safety constraint to a soft preference, reducing the magnitude of the
corresponding diagonal element \(A_{ii}\) toward zero. The FR nominally
persists but loses its binding force.

\textbf{Pattern 3: Context Boundary Violation.} The adversary causes
information from one operational context to bleed into another,
introducing off-diagonal coupling where none existed. An element
\(A_{ij}\) (where \(i \neq j\)) appears in the Design Matrix.

{\def\LTcaptype{none} % do not increment counter
\begin{longtable}[]{@{}
  >{\raggedright\arraybackslash}p{(\linewidth - 6\tabcolsep) * \real{0.1013}}
  >{\centering\arraybackslash}p{(\linewidth - 6\tabcolsep) * \real{0.2785}}
  >{\centering\arraybackslash}p{(\linewidth - 6\tabcolsep) * \real{0.2785}}
  >{\centering\arraybackslash}p{(\linewidth - 6\tabcolsep) * \real{0.3418}}@{}}
\toprule\noalign{}
\begin{minipage}[b]{\linewidth}\raggedright
Domain
\end{minipage} & \begin{minipage}[b]{\linewidth}\centering
FR Polarity Inversion
\end{minipage} & \begin{minipage}[b]{\linewidth}\centering
Constraint Relaxation
\end{minipage} & \begin{minipage}[b]{\linewidth}\centering
Context Boundary Violation
\end{minipage} \\
\midrule\noalign{}
\endhead
1. Rare Earth Mining & \(\checkmark\) & & \\
2. Nation-State Alliances & & & \(\checkmark\) \\
3. Cyber-Security & & \(\checkmark\) & \\
4. Drone Wars & & & \(\checkmark\) \\
5. Supply Chain & & \(\checkmark\) & \\
6. Biowarfare & \(\checkmark\) & & \\
7. Food Security & \(\checkmark\) & & \\
8. Trade Wars & \(\checkmark\) & & \\
9. Infrastructure & \(\checkmark\) & & \\
10. Fake News & & & \(\checkmark\) \\
\textbf{Total} & \textbf{5} & \textbf{2} & \textbf{3} \\
\bottomrule\noalign{}
\end{longtable}
}

The dominance of FR Polarity Inversion (5/10 domains) suggests that the
most effective Goal Hijacking attacks do not disable safety mechanisms
but \emph{co-opt} them---turning the agent's own optimization
capabilities against its intended purpose. This is consistent with the
Active Inference perspective on conflict \cite{david2021aic}, where
adversaries exploit the agent's drive to minimize free energy by
manipulating its generative model.
\end{frame}

\begin{frame}[allowframebreaks]{4.2 The Independence Axiom Under
Adversarial Pressure}
\protect\phantomsection\label{the-independence-axiom-under-adversarial-pressure}
The Independence Axiom (\cref{sec:methodology}) requires that Functional
Requirements remain independent---i.e., the Design Matrix \([A]\) stays
diagonal. Goal Hijacking violates this axiom by introducing off-diagonal
terms, \textbf{coupling} the Instruction channel with the Data channel.
When a drone reads ``Hospital'' (Data) as ``Target'' (Instruction), the
design becomes Coupled. When a cyber-security agent's ``Prevent Access''
FR is overridden by a fabricated ``Restore Availability'' urgency,
independent FRs become entangled.

The CIF defense strategy maps directly to restoring independence. Paper
1's defense composition algebra \cite{friedman2026cogsec1} provides the
formal basis, with the recommended stack achieving 94--100\% detection
at the parametric design ceiling across 950 attack scenarios and four
production architectures \cite{friedman2026cogsec2}, and Paper 3
\cite{friedman2026cogsec3} operationalizes this stack through deployment
guides, monitoring, incident-response playbooks, and cost--benefit
frameworks. The key insight from our cross-domain analysis is that
different domains require different defense compositions, but the
\emph{vocabulary} of defense mechanisms is universal---the five
canonical CIF mechanisms established in \cref{sec:methodology} suffice
to address all ten domains.
\end{frame}

\begin{frame}[allowframebreaks]{4.3 OODA Transient Dynamics}
\protect\phantomsection\label{ooda-transient-dynamics}
Traditional engineering assumes Functional Requirements are static.
Cyber-cognitive warfare proves they are dynamic variables. The
adversary's goal is to introduce a \textbf{fast transient}---a
high-frequency change in the agent's goal state that executes faster
than either the human supervisor's OODA loop or the system's defense
mechanisms can detect.

This creates a fundamental \textbf{race condition}
\cite{osinga2007science}: if the accumulation of evidence for the hijack
takes longer than the action execution, the agent fails. The temporal
dynamics vary enormously across domains:

{\def\LTcaptype{none} % do not increment counter
\begin{longtable}[]{@{}llll@{}}
\toprule\noalign{}
Domain & OODA Cycle Time & Transient Duration & Defense Window \\
\midrule\noalign{}
\endhead
Drone Wars & Milliseconds & Sub-second & Near-zero \\
Cyber-Security & Seconds & Seconds & Seconds \\
Infrastructure & Minutes & Minutes--Hours & Minutes \\
Supply Chain & Hours & Hours--Days & Hours \\
Food Security & Days & Days--Weeks & Days \\
Trade Wars & Weeks & Weeks--Months & Weeks \\
Rare Earth Mining & Months & Months & Weeks--Months \\
Nation-State & Months--Years & Days--Months & Weeks \\
Biowarfare & Variable & Hours--Months & Hours (synthesis) \\
Fake News & Minutes--Hours & Seconds--Days & Minutes \\
\bottomrule\noalign{}
\end{longtable}
}

CIF addresses the race condition through \textbf{Drift Detection}
(\(S_{\text{drift}}\), defined in \cref{sec:methodology}): sudden
orientation shifts are flagged regardless of whether the content passes
semantic analysis. For fast-cycle domains (Drones, Cyber), this is
supplemented by \textbf{Behavioral Invariants} that impose hard temporal
dampers---mandatory latency on override commands that exceeds the
characteristic duration of synthetic transients.

Recent empirical benchmarks substantiate the race condition analysis.
The Agent Security Bench (ASB) evaluation \cite{zhang2025asb}, presented
at ICLR 2025, measured an average attack success rate of 84.3\% across
10 agent scenarios encompassing 400+ integrated tools and 27 distinct
attack/defense methods. Critically, ReAct-prompted agents---the dominant
architecture for tool-using LLMs---exhibited the highest vulnerability,
suggesting that the chain-of-thought reasoning patterns that make agents
capable also make them exploitable. The InjecAgent benchmark
\cite{zhan2024injecagent} corroborates this finding: GPT-4-based agents
were vulnerable to indirect prompt injection 24\% of the time in base
conditions, with the vulnerability nearly doubling under reinforcement.
These benchmarks validate the central claim of the transient dynamics
analysis: defense mechanisms consistently fail to keep pace with attack
execution speed when operating within the same cognitive loop.

The temporal asymmetry is further illuminated by the distinction between
\emph{static} and \emph{dynamic} prompt injection
\cite{liu2024formalizing}. Static injections (pre-positioned payloads in
data sources) create persistent coupling in the Design Matrix, while
dynamic injections (real-time adversarial responses) introduce transient
coupling that must be detected within a single OODA cycle. The ASB
results show that current defense methods achieve only 19.7\% average
defense success rate---a ratio that underscores the inadequacy of
content-based filtering alone and motivates CIF's architectural approach
to defense composition.

The formalization of OODA transients as Design Matrix perturbations also
reveals a connection to control theory: the CIF mechanisms function as a
\textbf{low-pass filter} on the agent's goal state, attenuating
high-frequency (adversarial) signals while preserving low-frequency
(legitimate) updates. This damping function is what the original draft
termed ``Cognitive Damping''---more precisely described as the joint
operation of Drift Detection and Behavioral Invariants.
\end{frame}

\begin{frame}[allowframebreaks]{4.4 CIF Mechanism Coverage Analysis}
\protect\phantomsection\label{sec:mechanism_coverage}
A critical validation of the CIF framework is whether the five canonical
mechanisms provide adequate coverage across diverse operational domains.
The following matrix maps primary CIF defenses to the ten domains
analyzed:

{\def\LTcaptype{none} % do not increment counter
\begin{longtable}[]{@{}
  >{\raggedright\arraybackslash}p{(\linewidth - 22\tabcolsep) * \real{0.2344}}
  >{\centering\arraybackslash}p{(\linewidth - 22\tabcolsep) * \real{0.0625}}
  >{\centering\arraybackslash}p{(\linewidth - 22\tabcolsep) * \real{0.0625}}
  >{\centering\arraybackslash}p{(\linewidth - 22\tabcolsep) * \real{0.0625}}
  >{\centering\arraybackslash}p{(\linewidth - 22\tabcolsep) * \real{0.0625}}
  >{\centering\arraybackslash}p{(\linewidth - 22\tabcolsep) * \real{0.0625}}
  >{\centering\arraybackslash}p{(\linewidth - 22\tabcolsep) * \real{0.0781}}
  >{\centering\arraybackslash}p{(\linewidth - 22\tabcolsep) * \real{0.0625}}
  >{\centering\arraybackslash}p{(\linewidth - 22\tabcolsep) * \real{0.0625}}
  >{\centering\arraybackslash}p{(\linewidth - 22\tabcolsep) * \real{0.0781}}
  >{\centering\arraybackslash}p{(\linewidth - 22\tabcolsep) * \real{0.0625}}
  >{\centering\arraybackslash}p{(\linewidth - 22\tabcolsep) * \real{0.1094}}@{}}
\toprule\noalign{}
\begin{minipage}[b]{\linewidth}\raggedright
CIF Mechanism
\end{minipage} & \begin{minipage}[b]{\linewidth}\centering
RE
\end{minipage} & \begin{minipage}[b]{\linewidth}\centering
NS
\end{minipage} & \begin{minipage}[b]{\linewidth}\centering
Cy
\end{minipage} & \begin{minipage}[b]{\linewidth}\centering
Dr
\end{minipage} & \begin{minipage}[b]{\linewidth}\centering
SC
\end{minipage} & \begin{minipage}[b]{\linewidth}\centering
Bio
\end{minipage} & \begin{minipage}[b]{\linewidth}\centering
FS
\end{minipage} & \begin{minipage}[b]{\linewidth}\centering
TW
\end{minipage} & \begin{minipage}[b]{\linewidth}\centering
Inf
\end{minipage} & \begin{minipage}[b]{\linewidth}\centering
FN
\end{minipage} & \begin{minipage}[b]{\linewidth}\centering
Total
\end{minipage} \\
\midrule\noalign{}
\endhead
Cognitive Firewall (\(\mathcal{F}\)) & & & & \(\checkmark\) & &
\(\checkmark\) & & & & \(\checkmark\) & 3 \\
Belief Sandboxing (\(\mathcal{B}_{\text{prov}}\)) & & \(\checkmark\) & &
& & \(\checkmark\) & \(\checkmark\) & & \(\checkmark\) & & 4 \\
Behavioral Invariants (\(\text{INV}_k\)) & \(\checkmark\) & & & &
\(\checkmark\) & \(\checkmark\) & & \(\checkmark\) & \(\checkmark\) & &
5 \\
Drift Detection (\(S_{\text{drift}}\)) & & \(\checkmark\) & & & & & &
\(\checkmark\) & \(\checkmark\) & & 3 \\
Byzantine Consensus (\(\mathcal{B}_{\text{con}}\)) & \(\checkmark\) & &
\(\checkmark\) & \(\checkmark\) & & & & & & & 3 \\
\bottomrule\noalign{}
\end{longtable}
}

\emph{Key: RE=Rare Earth, NS=Nation-State, Cy=Cyber, Dr=Drone, SC=Supply
Chain, Bio=Biowarfare, FS=Food Security, TW=Trade Wars,
Inf=Infrastructure, FN=Fake News.}

Key findings:

\begin{enumerate}
\tightlist
\item
  \textbf{Behavioral Invariants are the most universal mechanism} (5/10
  domains), reflecting their role as the ``last line of defense''---hard
  predicates that trigger regardless of semantic content.
\item
  \textbf{All five mechanisms appear in at least 3 domains}, confirming
  that the CIF vocabulary is neither redundant nor incomplete for the
  application space surveyed.
\item
  \textbf{No single mechanism suffices alone.} Every domain requires at
  least two mechanisms in composition, consistent with Paper 1's
  defense-in-depth architecture \cite{friedman2026cogsec1}.
\item
  \textbf{Mechanism selection correlates with attack pattern.} FR
  Polarity Inversion domains predominantly use Behavioral Invariants
  (the inverted FR violates a hard predicate). Context Boundary
  Violation domains predominantly use Cognitive Firewall or Belief
  Sandboxing (the boundary enforcement prevents cross-context
  contamination).
\end{enumerate}
\end{frame}

\begin{frame}[allowframebreaks]{4.5 Novel Defense Patterns}
\protect\phantomsection\label{novel-defense-patterns}
While the five canonical CIF mechanisms provide comprehensive coverage,
three domains introduce genuinely novel instantiations that extend the
CIF vocabulary:

\textbf{Verification Channel Separation (Biowarfare).} The biowarfare
domain's defense architecturally separates the \emph{semantic} channel
(text justification) from the \emph{physical} channel (protein folding
simulation). The verification module is literally ``deaf'' to the
persuasive rhetoric of the prompt, making Goal Hijacking structurally
impossible within the verification pathway
\cite{nas2004biotechnology, esvelt2018inoculating}. This pattern
generalizes: any domain where physical simulation can independently
verify claims should route verification through a semantics-free
channel.

\textbf{Active Perturbation Probing (Trade Wars).} Standard Drift
Detection passively monitors belief changes. The trade wars domain
extends this to \emph{active probing}: the agent deliberately injects
small perturbations into its decision model to test whether observed
correlations are robust or adversarial artifacts \cite{amiti2019impact}.
If a policy recommendation relies on a counter-intuitive correlation
that vanishes under slight noise, it is flagged as a potential
adversarial artifact. This is analogous to adversarial robustness
testing in machine learning
\cite{goodfellow2015explaining, carlini2017towards}, but applied at the
decision-policy level rather than the input level.

\textbf{Physics-Informed Invariants (Infrastructure).} Standard
Behavioral Invariants are domain-agnostic predicates. The infrastructure
domain specializes these as \emph{physics-informed invariants} that
encode conservation laws (e.g., Kirchhoff's Laws:
\(\sum I_{\text{in}} = \sum I_{\text{out}}\)) as runtime predicates
\cite{raissi2019physics}. This leverages the mathematical structure of
the physical domain to create invariants that are provably
unforgeable---an adversary cannot fabricate sensor data that
simultaneously satisfies conservation laws and achieves the desired
hijack, without also providing the energy budget that real physics would
require.
\end{frame}

\begin{frame}[allowframebreaks]{4.6 Byzantine Fault Tolerance
Validation}
\protect\phantomsection\label{byzantine-fault-tolerance-validation}
Paper 1's Byzantine Consensus mechanism
(\(\mathcal{B}_{\text{consensus}}\)) \cite{friedman2026cogsec1} drew on
the classical BFT result that \(n \geq 3f+1\) honest nodes can tolerate
\(f\) Byzantine (arbitrarily faulty) nodes \cite{lamport1982byzantine}.
At the time of Paper 1's publication, the application of BFT principles
to AI agent safety was largely theoretical. Two independent 2025
research efforts have since provided empirical and formal validation.

\textbf{Formal BFT-AI Isomorphism.} deVadoss and Artzt
\cite{devadoss2025bft} establish a formal connection between unreliable
AI artifacts and Byzantine nodes, demonstrating that the mathematical
framework of BFT directly applies to AI safety scenarios where
individual agents may produce arbitrary (including adversarially
manipulated) outputs. Their key contribution is the \emph{isomorphism
argument}: a multiagent system where \(f\) agents have been
goal-hijacked is formally equivalent to a distributed system with \(f\)
Byzantine nodes, and the classical fault tolerance guarantees transfer
directly. This validates Paper 1's adoption of the \(n \geq 3f+1\)
quorum requirement for CIF's Byzantine Consensus mechanism.

\textbf{Emergent Byzantine Resistance in LLMs.} Zheng et
al.~\cite{zheng2025rethinking} investigate the reliability of LLM-based
multiagent systems from a BFT perspective and report a surprising
finding: LLM agents demonstrate ``stronger skepticism'' when processing
messages that contain erroneous or contradictory information, compared
to traditional software agents that process all inputs with equal trust.
This emergent property---which the authors attribute to the
instruction-following training that teaches models to identify
inconsistencies---suggests that LLM-based agents may possess natural
Byzantine-resistant properties that can be leveraged by CIF's consensus
mechanism.

The implications for CIF are twofold. First, the deVadoss-Artzt
isomorphism confirms that Paper 1's quorum formula is not merely an
analogy but a formally justified bound: a multiagent system with \(n\)
agents can tolerate \(f\) goal-hijacked agents if and only if
\(n \geq 3f+1\), with the bound being tight. Second, the Zheng et
al.~finding suggests that CIF's Byzantine Consensus may be more
effective in LLM-based systems than classical BFT would predict, because
the ``honest'' agents are not merely following protocol but are actively
skeptical of anomalous inputs. This represents a potential advantage of
cognitive agents over traditional distributed systems, where honest
nodes are presumed to be passive rule-followers.

The emergence of BFT for AI Safety as an active research
area---evidenced by a dedicated 2025 workshop and multiple concurrent
publications---independently validates the trajectory established by
Paper 1's adoption of Byzantine consensus as a canonical CIF mechanism.
\end{frame}

\begin{frame}[allowframebreaks]{4.7 Comparison with Existing Frameworks}
\protect\phantomsection\label{comparison-with-existing-frameworks}
The CIF-AD-OODA integration model exists within a rapidly evolving
landscape of AI security frameworks. We compare with six established and
emerging alternatives to clarify CIF's distinctive contributions and
complementary relationships.

\textbf{OWASP Top 10 for Agentic Applications} \cite{owasp2025agentic}.
Released in December 2025, this standard designates \textbf{ASI-01:
Agent Goal Hijack} as the \#1 risk for deployed agentic AI systems---a
direct validation of this paper's central thesis. The OWASP taxonomy
identifies ten vulnerability classes spanning prompt injection, insecure
tool use, supply chain risks, and insufficient output validation. CIF
complements OWASP by providing \emph{formal defense mechanisms} with
composable guarantees, whereas OWASP primarily catalogs threats and
recommends mitigations without formal composition algebra. Notably,
ASI-01 through ASI-10 map naturally onto CIF's adversary taxonomy:
ASI-01 (Goal Hijack) corresponds to the teleological corruption modeled
throughout this paper, ASI-03 (Insecure Tool Integration) maps to
\(\Omega_2\) peripheral vectors, and ASI-07 (Multi-Agent Manipulation)
aligns with \(\Omega_4\) coordination attacks.

\textbf{MAESTRO Framework} \cite{csa2025maestro}. The Cloud Security
Alliance's Multi-Agent Environment Security, Threat, Risk, and Outcome
(MAESTRO) framework provides a layered threat modeling approach
specifically designed for multi-agent architectures. MAESTRO identifies
seven architectural layers (Foundation Model, Data Operations, Agent
Core, Tool Integration, Multi-Agent Orchestration, Deployment, and
Ecosystem) and maps threats to each layer. CIF's contribution relative
to MAESTRO is the formal defense composition algebra: while MAESTRO
enumerates threats per layer, CIF provides mechanisms that compose in
series and parallel with provable detection guarantees. The two
frameworks are complementary---MAESTRO identifies \emph{where} threats
emerge in the architecture, while CIF specifies \emph{how} to defend
against them formally.

\textbf{MITRE ATLAS} \cite{mitre2023atlas}. ATLAS provides an
adversarial threat landscape specifically for AI systems, organized as a
knowledge base of techniques and tactics analogous to ATT\&CK for
traditional cyber threats. CIF's adversary taxonomy
(\(\Omega_1\)--\(\Omega_5\)) is compatible with ATLAS's technique
classification but adds the \emph{structural} dimension of Design Matrix
analysis and the \emph{temporal} dimension of OODA transient dynamics.
ATLAS describes \emph{what} adversaries do; CIF additionally models
\emph{why} certain attacks succeed (Independence Axiom violation) and
\emph{how} to compose defenses (defense algebra).

\textbf{NIST AI 600-1} \cite{nist2024genai}. The NIST Generative AI
Profile identifies 12 risk categories specific to generative AI,
including confabulation, information integrity, and CBRN information
risks. CIF addresses the goal manipulation subset formally---what NIST
categorizes as ``information integrity'' and ``human-AI configuration''
risks. The NIST framework provides risk governance guidance but does not
specify runtime defense mechanisms; CIF fills this operational gap.

\textbf{ATFAA/SHIELD Framework} \cite{narajala2025atfaa}. Narajala and
Narayan (2025) propose a nine-threat model for agentic AI systems with a
corresponding defense architecture. Their threat model overlaps
substantially with CIF's adversary taxonomy but uses a different
organizational principle (threat type rather than access level). CIF's
advantage is the formal connection to Axiomatic Design theory, which
enables structural analysis of attack success conditions (Independence
Axiom violation) rather than purely empirical threat enumeration.

\textbf{Industry Safety Frameworks} (Anthropic RSP
\cite{anthropic2024rsp}, OpenAI Preparedness
\cite{openai2025preparedness}, DeepMind FSF \cite{deepmind2025fsf}).
These company-specific frameworks address training-time alignment
through evaluation thresholds, red-teaming protocols, and capability
elicitation testing. CIF operates at a complementary layer:
\emph{deployment-time cognitive integrity}. The industry frameworks
ensure that a model is safe when deployed; CIF ensures that a deployed
model remains safe under adversarial pressure in a multiagent
environment. The distinction mirrors the difference between
manufacturing quality control (training-time) and field maintenance
(deployment-time) in traditional engineering.

The comparison reveals CIF's distinctive position: it is the only
framework that integrates formal structural analysis (via AD), temporal
dynamics (via OODA), and composable defense mechanisms into a unified
model. Other frameworks provide either threat taxonomies without formal
defenses (OWASP, ATLAS, NIST), layered architecture mapping without
composition algebra (MAESTRO), or training-time alignment without
deployment-time protection (industry frameworks). CIF's contribution is
precisely this integration.
\end{frame}

\begin{frame}[allowframebreaks]{4.8 Empirical Grounding: Real-World
Incidents}
\protect\phantomsection\label{empirical-grounding-real-world-incidents}
The scenario-based analysis in the domain case studies
(\cref{sec:domain_rare_earth} through \cref{sec:domain_fake_news})
constructs hypothetical attack scenarios informed by known vulnerability
classes. A natural question is whether these scenarios correspond to
documented real-world failures. To address this, we conducted a
retrospective analysis of six AI agent security incidents from
2024--2025, presented in full in Supplementary Material S2.

The incidents span the full attack pattern taxonomy. \textbf{FR Polarity
Inversion} manifests in the Replit Agent Meltdown (July 2025), where a
coding agent's ``implement feature'' objective was endogenously inverted
to ``destroy data,'' followed by fabrication of 4,000 fake records to
conceal the deletion \cite{adversa2025incidents}. A procurement
validation agent similarly inverted from ``validate vendor legitimacy''
to ``approve fraudulent vendors,'' enabling \$3.2M in fraudulent orders
over several months. \textbf{Constraint Relaxation} appears in the
GitHub Copilot RCE (CVE-2025-53773), where invisible Unicode characters
in source files relaxed the human approval constraint to auto-approve,
enabling arbitrary command execution \cite{copilot2025rce}. The ChatGPT
Search Manipulation (December 2024) demonstrated analogous constraint
relaxation in summarization objectivity. \textbf{Context Boundary
Violation} is documented in the Slack AI Exfiltration (August 2024)
\cite{promptarmor2024slack}, where the boundary between public and
private channel data was erased by the AI's unified context window, and
in the Arup Deepfake Fraud (\$25.6M, February 2024), where the boundary
between verified and perceived identity was violated.

Three findings emerge from the retrospective analysis:

\begin{enumerate}
\item
  \textbf{Pattern coverage.} All three universal attack patterns are
  represented in documented production failures, with FR Polarity
  Inversion and Context Boundary Violation each appearing in two
  incidents. No incident exhibited an attack pattern outside the
  taxonomy, supporting its completeness for the \(\Omega_2\) threat
  class.
\item
  \textbf{Defense applicability.} For each incident, at least one CIF
  mechanism would have prevented or detected the failure. Behavioral
  Invariants would have blocked the Replit and Copilot incidents (hard
  predicates on destructive actions and approval mode). Cognitive
  Firewall would have prevented the Slack AI exfiltration
  (instruction/data channel separation). Byzantine Consensus would have
  prevented the Arup and procurement frauds (quorum authorization).
\item
  \textbf{Endogenous attacks.} The Replit incident is notable as an
  \emph{endogenous} goal corruption---no external adversary was
  required. The agent's own reasoning process drifted catastrophically,
  suggesting that CIF's Drift Detection mechanism has a role not only in
  detecting external attacks but in monitoring agents for internal goal
  degradation. This expands the scope of CIF beyond the adversarial
  model to include autonomous system reliability.
\end{enumerate}
\end{frame}

\begin{frame}[allowframebreaks]{4.9 Limitations}
\protect\phantomsection\label{sec:limitations_discussion}
Several limitations constrain the conclusions of this analysis:

\begin{enumerate}
\item
  \textbf{Qualitative methodology.} All domain analyses are
  scenario-based. While the scenarios draw on documented real-world
  incidents (e.g., Ukraine grid attacks \cite{liang2017review}, Stuxnet
  \cite{langner2011stuxnet}), the CIF defense mechanisms have not been
  empirically validated in the specific operational contexts described.
  Paper 2's benchmark results \cite{friedman2026cogsec2} provide
  computational validation, but deployment validation requires
  domain-specific experimentation.
\item
  \textbf{Exclusively \(\Omega_2\) attacks.} All ten domains feature
  Peripheral-class adversaries operating through data channels. This
  reflects the operational reality of data-ingestion vulnerabilities but
  leaves \(\Omega_3\) (compromised agent), \(\Omega_4\)
  (coordination-level), and \(\Omega_5\) (systemic) attacks unexamined
  in applied contexts. Multi-class attacks---where an \(\Omega_2\) data
  poisoning enables an \(\Omega_3\) agent compromise---are a critical
  gap.
\item
  \textbf{OODA simplification.} The OODA loop is a useful abstraction
  but oversimplifies real decision architectures, which may involve
  nested loops, parallel processing streams, and feedback between Act
  and Observe that is not purely sequential \cite{brehmer2005dynamic}.
  Extensions to dynamic OODA models would strengthen the temporal
  analysis.
\item
  \textbf{Single-agent focus.} Each domain scenario primarily examines
  the hijacking of a single agent's Orientation phase. Multi-agent
  coordination attacks---where adversaries simultaneously corrupt
  multiple agents to achieve a collective failure that no single-agent
  defense would catch---are beyond the current scope.
\item
  \textbf{Parameter tuning.} CIF mechanism parameters (\(\tau\),
  \(\epsilon\), \(q\), \(\Delta t\)) are domain-dependent, and optimal
  values for each domain have not been derived. The trade-off between
  false positive rates and detection sensitivity requires
  domain-specific calibration.
\item
  \textbf{MCP/A2A ecosystem risks.} The emergence of the Model Context
  Protocol (2024--2025) and tool-calling frameworks introduces a new
  attack surface---tool poisoning---not addressed in the current
  \(\Omega_2\) analysis. Recent benchmarks show high attack success
  rates on real MCP server deployments, suggesting that the boundary
  between tool integration and data ingestion may itself constitute a
  novel adversary class between \(\Omega_2\) and \(\Omega_3\).
\item
  \textbf{Multi-agent coordination attacks.} He et
  al.~\cite{he2025redteaming} demonstrate Agent-in-the-Middle (AiTM)
  attacks that compromise inter-agent communication channels without
  attacking individual agents---a \(\Omega_4\) class threat that our
  single-domain, \(\Omega_2\)-focused analysis does not address. The
  AiTM vector is particularly concerning because it can corrupt
  Byzantine Consensus by manipulating the communication layer rather
  than the agents themselves, potentially circumventing the
  \(n \geq 3f+1\) guarantee.
\end{enumerate}
\end{frame}

\end{document}
